% Generated draft: requires mathematical and source review.
\documentclass[11pt]{article}
\usepackage[margin=25mm]{geometry}
\usepackage{fontspec}
\setmainfont{Latin Modern Roman}
\usepackage{amsmath,amssymb,mathtools,unicode-math}
\setmathfont{Latin Modern Math}
\usepackage{enumitem,needspace,xurl,hyperref}
\setlength{\parindent}{0pt}
\setlength{\parskip}{5pt}
\setlength{\emergencystretch}{3em}
\newcommand{\N}{\mathbb{N}}
\newcommand{\Z}{\mathbb{Z}}
\newcommand{\Q}{\mathbb{Q}}
\newcommand{\R}{\mathbb{R}}
\newcommand{\C}{\mathbb{C}}
\newcommand{\F}{\mathbb{F}}
\newcommand{\A}{\mathbb{A}}
\newcommand{\PP}{\mathbb{P}}
\newcommand{\E}{\mathbb{E}}
\newcommand{\eps}{\varepsilon}
\newcommand{\dd}{\,\mathrm{d}}
\begin{document}
\section*{UnsolvedMath ID: 30002701}
\textbf{Source problem code:} catalogue entry
\subsection{Definitions and mathematical statement}
The supplied problem concerns an unspecified dynamical system with map \(f\), observables \(\varphi_1,\ldots,\varphi_k\), and objects denoted by \(W^s\) and \(W^u\), intended as stable and unstable foliations. The underlying space, measure, hypotheses on \(f\), admissible observables, and precise definitions of these foliations are not supplied. Consequently, the following is a conditional formulation of the recorded question, with these data left unresolved.

Write \((X,\mu)\) for the unspecified measured space underlying \(L^2\). Interpreting the iteration gaps as positive integers, for a fixed finite number \(k\geq 1\) of observables and \(\mathbf n=(n_1,\ldots,n_k)\), put
\[
S_j(\mathbf n)=\sum_{i=1}^{j}n_i,
\qquad
P_{\mathbf n}=\prod_{j=1}^{k}\bigl(\varphi_j\circ f^{S_j(\mathbf n)}\bigr).
\]
These expressions require the relevant iterates and products to be defined and the products to belong to \(L^2(X,\mu)\); the supplied record does not give hypotheses guaranteeing this. The range of \(k\) and the admissible class of observables also remain to be specified.

A weak-\(L^2\) limit along gaps tending to infinity means a function \(g\in L^2(X,\mu)\) and a sequence \(\mathbf n^{(r)}\) such that
\[
\lim_{r\to\infty}n_i^{(r)}=+\infty
\quad(1\leq i\leq k),
\qquad
P_{\mathbf n^{(r)}}\rightharpoonup g
\quad\text{in }L^2(X,\mu).
\]
Explicitly, the weak convergence condition is
\[
\lim_{r\to\infty}\int_X P_{\mathbf n^{(r)}}\,\overline h\,\dd\mu
=
\int_X g\,\overline h\,\dd\mu
\qquad\text{for every }h\in L^2(X,\mu),
\]
with complex conjugation omitted in a real-valued formulation. The source does not specify whether the observables are real or complex.

\textbf{Recorded question.} Given the asserted \(W^s\)-invariance of every such weak-\(L^2\) accumulation point, must every such accumulation point also be \(W^u\)-invariant?

Here foliation invariance is intended to express constancy along the corresponding leaves, but its precise almost-everywhere or representative-based meaning must be recovered from the source. Since an \(L^2\) function is an equivalence class modulo null sets, leafwise constancy requires a specified convention. The question does not assert the existence of an accumulation point. The missing dynamical hypotheses and invariance convention prevent a fully specified universal mathematical statement.
\subsection{Short English statement}
The question asks whether every weak-\(L^2\) limit of the displayed products, when every successive iteration gap tends to infinity, is invariant along unstable leaves, given the recorded stable invariance. The dynamical setting, admissible observables, and precise meaning of leafwise invariance are missing and require mathematical review.
\subsection{Sources}
\begin{itemize}
\item Local UnsolvedMath registry: \texttt{data/unsolvedmath/problems.json}, permanent ID \(30002701\), title ``Unstable Invariance of Multiple-Correlation Limits,'' exact source problem code \texttt{catalogue entry}, in the Oberwolfach Reports Open Problems collection. No individual proposer is supplied; the proposed year is recorded as 2014.
\item Supplied source citation: \emph{Mini-Workshop: Differentiable Ergodic Theory, Dimension Theory and Stable Foliations} (2015), Oberwolfach Reports, pp.~2569--2617, DOI \texttt{10.4171/owr/2014/46}. \url{https://doi.org/10.4171/owr/2014/46}. The literature-source list labels this reference ``Oberwolfach Report 46/2014.'' The differing year labels are retained without resolution.
\item Supplied literature reference: Coudene, \emph{On Invariant Distributions and Mixing}. \url{https://doi.org/10.1017/S0143385706000733}.
\item Supplied literature reference: \emph{Multiple Mixing from Weak Hyperbolicity}. \url{https://arxiv.org/abs/1406.4015}.
\item Literature assessment date: \texttt{2026-08-21}. This is supplied dated curation, not independently verified. The registry labels the problem open and reports that stable invariance is proved, unstable invariance is explicitly posed as open, and later weak-hyperbolicity results establish mixing under additional assumptions rather than a general resolution. These are supplied source assessments, not independently established conclusions of this draft. The background explicitly requests independent verification of time-sensitive status and cited claims.
\item The supplied registry labels the statement \texttt{corrected\_verified} and reports a manual semantic, typographic, and LaTeX review. That supplied review label does not resolve the missing mathematical hypotheses identified here. No external source was visited for this draft.
\end{itemize}
\end{document}
